%%% Local Variables:
%%% mode: LaTeX
%%% TeX-engine: luatex
%%% TeX-master: t
%%% TeX-backend: "biber"
%%% End:
\documentclass[article,spanish]{apuntes-scr}

\subject{Manual de Referencia Oficial}
\title{Guía Completa de Uso y Manejo de Sway y Waybar}
\subtitle{Gestor de Ventanas en Mosaico (Tiling) y Barra de Estado para Wayland}
\author{Campos Mendoza, D. Daniel}
\date{%
    \textit{Ecosistema GNU Emacs 30 / Sway / Wayland} \\[0.3em]
    \today\ \textperiodcentered\ \DTMcurrenttime
}

% =========================================================
% BIBLIOGRAFÍA
% =========================================================
% \addbibresource{references.bib}

\begin{document}

\maketitle
\tableofcontents
\vspace{1.5em}\hrule\vspace{1.5em}

\section{Arquitectura del Sistema Sway}
\textbf{Sway} es un compositor y gestor de ventanas en mosaico (\textit{Tiling Window Manager}) nativo para Wayland, 100\% compatible con la sintaxis de i3 pero optimizado para el servidor gráfico moderno de Linux sin dependencia de X11/XWayland. Opera bajo el principio de teclado-primero (\textit{keyboard-centric}) eliminando la necesidad del ratón.

El modificador maestro del sistema es la tecla \textbf{Super / Windows}:
$$\mathbf{\$mod = \text{Mod4}}$$

\section{Atajos de Teclado del Ecosistema}

\subsection{Lanzadores Principales y Control de Sesión}
\begin{itemize}[leftmargin=*]
    \item \textbf{\texttt{\$mod + Return}}: Abre el emulador de terminal \textbf{Kitty}.
    \item \textbf{\texttt{\$mod + d}}: Lanza el menú interactivo \textbf{Wofi} para buscar aplicaciones.
    \item \textbf{\texttt{\$mod + b}}: Lanza el navegador web \textbf{Firefox} (dirigido automáticamente al Workspace 4).
    \item \textbf{\texttt{\$mod + q}}: Cierra de forma ordenada la ventana activa (\textit{kill}).
    \item \textbf{\texttt{\$mod + Shift + c}}: Recarga la configuración de Sway en caliente sin cerrar ninguna aplicación.
    \item \textbf{\texttt{\$mod + Shift + e}}: Cierra la sesión de Sway y regresa al gestor de acceso.
\end{itemize}

\subsection{Navegación y Foco de Ventanas (Estilo Vim)}
\begin{itemize}[leftmargin=*]
    \item \textbf{\texttt{\$mod + h}} o \textbf{\texttt{\$mod + $\leftarrow$}}: Enfocar ventana a la \textbf{izquierda}.
    \item \textbf{\texttt{\$mod + j}} o \textbf{\texttt{\$mod + $\downarrow$}}: Enfocar ventana \textbf{inferior}.
    \item \textbf{\texttt{\$mod + k}} o \textbf{\texttt{\$mod + $\uparrow$}}: Enfocar ventana \textbf{superior}.
    \item \textbf{\texttt{\$mod + l}} o \textbf{\texttt{\$mod + $\rightarrow$}}: Enfocar ventana a la \textbf{derecha}.
\end{itemize}

\subsection{Mover y Reorganizar Ventanas}
\begin{itemize}[leftmargin=*]
    \item \textbf{\texttt{\$mod + Shift + h / j / k / l}}: Mueve la ventana activa en la dirección indicada.
\end{itemize}

\subsection{Distribuciones de Pantalla (Layouts)}
\begin{itemize}[leftmargin=*]
    \item \textbf{\texttt{\$mod + e}}: Alterna la dirección del split para la próxima ventana (\textit{toggle split}).
    \item \textbf{\texttt{\$mod + h}} / \textbf{\texttt{\$mod + v}}: Pre-fija división horizontal / vertical.
    \item \textbf{\texttt{\$mod + w}} (\textit{Tabbed}): Agrupa todas las ventanas en \textbf{pestañas superiores}.
    \item \textbf{\texttt{\$mod + s}} (\textit{Stacking}): Apila las ventanas mostrando sus barras de título verticalmente.
    \item \textbf{\texttt{\$mod + f}}: Alterna el modo \textbf{Pantalla Completa} (\textit{Fullscreen}).
    \item \textbf{\texttt{\$mod + Shift + Espacio}}: Alterna la ventana entre modo \textbf{Flotante} y modo \textbf{Mosaico}.
    \item \textbf{\texttt{\$mod + Espacio}}: Alterna el foco entre las ventanas flotantes y el mosaico.
\end{itemize}

\section{Organización de Espacios de Trabajo (Workspaces)}
El flujo de trabajo se estructura en 5 escritorios dedicados:
\begin{center}
\begin{tabular}{lll}
\toprule
\textbf{Workspace} & \textbf{Etiqueta} & \textbf{Propósito} \\
\midrule
\texttt{\$mod + 1} & \textbf{1: Term} & Terminales principales Kitty \\
\texttt{\$mod + 2} & \textbf{2: Code} & GNU Emacs 30 (Edición de código y \LaTeX{}) \\
\texttt{\$mod + 3} & \textbf{3: Doc} & Zathura (Lectura de libros, artículos y apuntes) \\
\texttt{\$mod + 4} & \textbf{4: Web} & Navegación e investigación con Firefox \\
\texttt{\$mod + 5} & \textbf{5: Media} & Mensajería, audio y reproducción multimedia \\
\bottomrule
\end{tabular}
\end{center}

\begin{itemize}[leftmargin=*]
    \item \textbf{\texttt{\$mod + [1 - 5]}}: Cambiar al Workspace indicado.
    \item \textbf{\texttt{\$mod + Shift + [1 - 5]}}: Mover la ventana activa al Workspace seleccionado.
\end{itemize}

\section{El Scratchpad (Terminal Flotante Emergente)}
\begin{itemize}[leftmargin=*]
    \item \textbf{\texttt{\$mod + \`{}} (Grave / Tilde)}: Abre o invoca a pantalla una terminal Kitty flotante y centrada (\texttt{900x600px}) para operaciones rápidas del sistema.
    \item \textbf{\texttt{\$mod + -}} / \textbf{\texttt{\$mod + Shift + -}}: Mostrar ventana del scratchpad / Enviar ventana activa al scratchpad.
\end{itemize}

\section{Utilidades de Hardware y Teclado}
\begin{itemize}[leftmargin=*]
    \item \textbf{Print (Impr Pant)}: Captura interactiva de un área de la pantalla (\texttt{grim + slurp}) copiada directamente al portapapeles.
    \item \textbf{Win + Espacio}: Alterna la distribución del teclado entre \textbf{Latinoamérica (\texttt{latam})}, \textbf{España (\texttt{es})} e \textbf{Inglés (\texttt{us})}.
    \item \textbf{Bloq Mayús (Caps Lock)}: Remapeada por hardware para actuar como tecla \textbf{Escape} para máxima ergonomía en modo Evil/Vim.
    \item \textbf{Control de Volumen y Brillo}: Teclas multimedia integradas con \texttt{wpctl} y \texttt{brightnessctl}.
\end{itemize}

\section{Barra de Estado Waybar}
Configurada en \texttt{\textasciitilde/.config/waybar/config} con estilo CSS oscuro:
Muestra en tiempo real los espacios de trabajo activos, uso porcentual de CPU y memoria RAM, nivel de batería, estado de volumen de audio y fecha/hora con tipografía \texttt{Iosevka Term}.

\end{document}
